\subsection{Actual endpoint scores and their intrinsic normalization}
\label{paper:ep-intrinsic-section}
We now distinguish the raw local cost from the cost that enters a native
wall.  This distinction matters when a platform inherits a large
entrance from its predecessor.
For a named endpoint star template $v$, let $A_{v,g}$ be its retained
exterior names in group $g$, $q_g=|A_{v,g}|$, and $l=l(v)$ its last
positive-rank group.  Its reference cost is
\[
 \overline C_v=\sum_gL_gg_g(q_g).
\]
Let $\mathcal B_v$ and $\mathcal G_v$ be respectively its bad and good
proper groups, as in Proposition~\ref{paper:ep-count-box}, and put
\[
 L_{\rm bad}=\sum_{g\in\mathcal B_v}L_g,\qquad
 L_{\rm zero}=T-V_l,\qquad
 S_{\rm good}(v)=\sum_{g\in\mathcal G_v}\sqrt{L_g}.
\]
All other groups are full and incur no exposed-name discrepancy.
The reference estimates and the displaced box give, simultaneously,
\begin{align}
 \overline C_v&\geq a_HL_{\rm bad},\qquad
 \overline C_v\geq b_H\min(V_l,T-V_l),\nonumber\\
 C_v&\geq b_0+Z_l+\tfrac34\overline C_v+G_v-C_{\rm box},
 \qquad G_v\geq c_HS_{\rm good}(v)-C_H.
 \label{paper:ep-corner-input}
\end{align}
The last comparison follows from own-scale centrality: on a large group
$N_g\geq c_HL_g$, whereas bounded $N_g$ forces bounded $L_g$.

\begin{lemma}[Early and late corner payment]\label{paper:ep-corner-payment}
On this single box, every named endpoint star corner obeys
\begin{align}
 V_l\leq T/2:&\qquad C_v\geq a_0+c_HV_l-C_H,
                       \label{paper:ep-early-payment}\\
 V_l\geq T/2:&\qquad C_v\geq b_0+c_HL_{\rm zero}
                 +c_HL_{\rm bad}+c_HS_{\rm good}(v)-C_H.
                       \label{paper:ep-late-payment}
\end{align}
The all-zero and complete-deadline corners have costs exactly $a_0$
and $b_0$, respectively.
\end{lemma}
\begin{proof}
In the early case use $b_0+Z_l=a_0-\sum_{g\leq l}\zeta_g$.
Its error is $O(\sqrt{1+V_l})$ and is absorbed once by the positive
prefix charge in \eqref{paper:ep-corner-input}, retaining a fixed
fraction of $V_l$.  In the late case use
$b_0+Z_l=b_0+\sum_{g>l}\zeta_g$; only the zero-rank tail has a
square-root error.  Split the $3\overline C_v/4$ budget into two
equal pieces for $L_{\rm bad}$ and $L_{\rm zero}$, absorb that error
with half the latter piece, and retain the good-group gain.
The extreme identities follow directly from
\eqref{paper:ep-full-stop}.  The same proof, with a fixed constant,
includes $T=0$.
\end{proof}

Here is the exact cancellation behind two-sided scores.  With $N_g(A)$
the exterior count in $A$, write
\[
 H_g=\log S_g\,n_g-\log z\,k_g,\qquad
 H_g(A)=\log(q_g+z)[k_g+N_g(A)]-\log z\,k_g.
\]
For an empty group $H_g(A)=0$; for a full group it is $H_g$ exactly.
The notation $\overline H_g$ below means the deterministic natural
reference potential, including unconditioned old routing.  Set
\[
 \nu_{vg}=[H_g-H_g(A_{v,g})]
                  -[\overline H_g-\overline H_g(A_{v,g})].
\]
Then $\nu_{vg}=0$ identically on every full group, and direct expansion
of $C_v=a_0+\sum_{g\leq l}[q_gL_g-H_g(A_{v,g})]$ gives
\begin{equation}\label{paper:ep-exact-residual}
 C_v=a_0-\sum_{g\leq l}\zeta_g+\overline C_v
                   +\sum_{g\ {\rm proper}}\nu_{vg}
     =b_0+\sum_{g>l}\zeta_g+\overline C_v
                   +\sum_{g\ {\rm proper}}\nu_{vg}.
\end{equation}
Every repaired named count is within $C_H\sqrt{1+L_g}$ of its own
natural group mean.  Therefore
\begin{equation}\label{paper:ep-residual-central}
 |\nu_{vg}|\leq C_H\sqrt{1+L_g},\qquad
 \left|\sum_{g\in J}\zeta_g\right|
                  \leq C_H\sqrt{1+\sum_{g\in J}L_g}
\end{equation}
for consecutive $J$.  In the late formula the discrepancies of large
full groups have canceled, rather than contributing a spurious
$\sqrt T$ error.

\begin{lemma}[Raw finite scores]\label{paper:ep-raw-scores}
Suppose $a_0,b_0\geq0$ lie in fixed comparable intervals of deterministic
scales $A_{\rm in},B_{\rm in}$, respectively, allowing scale zero for
an identically zero entrance.  For a fixed sufficiently large $h$, set
\[
 S_v^{\rm raw}=\begin{cases}
 h+A_{\rm in}+\overline C_v,&V_l\leq T/2,\\
 h+B_{\rm in}+\overline C_v+S_{\rm good}(v),&V_l>T/2.
 \end{cases}
\]
Then $c_HS_v^{\rm raw}\leq C_v+h\leq C_HS_v^{\rm raw}$
simultaneously on every admitted count type.
\end{lemma}
\begin{proof}
For an early corner all proper groups are in the prefix, and
\eqref{paper:ep-exact-residual}--\eqref{paper:ep-residual-central}
give
$C_v\leq a_0+\overline C_v+C_H\sqrt{1+V_l}$.
Since $V_l\ll_H\overline C_v$, this is
$a_0+C_H'\overline C_v+C_H'$.
For the lower bound use the first expression for $b_0+Z_l$ in
\eqref{paper:ep-corner-input}, discard $G_v\geq0$, and absorb its
square root with at most $\overline C_v/4$ to get
\begin{equation}\label{paper:ep-early-raw}
 C_v\geq a_0+\tfrac12\overline C_v-C_H.
\end{equation}
For a late corner, the stronger box row similarly gives
\begin{equation}\label{paper:ep-late-raw}
 C_v\geq b_0+\tfrac12\overline C_v
                         +c_HS_{\rm good}(v)-C_H.
\end{equation}
The tail formula gives an upper bound by
$b_0+\overline C_v$ plus the square roots of the zero tail, bad groups
and good groups.  The first two are bounded by $C_H(1+\overline C_v)$;
the last are $C_HS_{\rm good}+C_H$.  A fixed large $h$ completes both
comparisons, including the literal extreme rows.
\end{proof}
These raw scores describe $C_v$.  They are not the native reserve
scores, which must account for the normalization in the full cost.

For later use we spell out how the raw payment interacts with exact
competitor classes.  Let $A$ be a finite nonnegative complete-type weight
in the own-group central support, $W=\mathbb E_\pi A$, and let
$\mathcal E_v=\mathcal B_v\cup\mathcal G_v\cup\{g>l(v)\}$
be its exposed groups.  Full groups have one cell containing all eligible
names.  The exact count-reveal estimate \eqref{ld:reveal-bound} therefore
gives, in the original class of the retained anchor,
\begin{equation}\label{paper:ep-corner-reveal}
 W_v:=\mathbb E_\pi[A\mid\operatorname{Class}_v]
 \leq C_HW\prod_{g\in\mathcal E_v}(1+n_g)^{(a_g-1)/2}.
\end{equation}
There is no full-group count factor and no class-probability factor.
If a predeclared selected-bridge cost satisfies
$D_g\geq c_D(1+n_g)^{-A_D}$ (and $D_g=1$ on unselected groups),
then \eqref{paper:ep-late-payment} pays every polynomial in
\eqref{paper:ep-corner-reveal} on a late corner, giving
\begin{equation}\label{paper:ep-late-class-payment}
 e^{-C_v}W_v\leq C_He^{-b_0}W\prod_{g\in\mathcal E_v}D_g.
\end{equation}
Indeed $n_g\ll_H1+L_g$, and a positive fraction of $L_g$ on a bad or
zero-tail group, or of $\sqrt{L_g}$ on a good group, dominates any
fixed multiple of $\log(1+L_g)$ up to a fixed constant.  The same
argument works for any fixed positive fraction of $C_v-b_0$.
Under the same-class bank hypotheses of \eqref{ld:class-banks}, apply
the bank theorem with every $\mathcal E_v$ slot exposed and depth zero
elsewhere.  Its bounds $q_{g,0}\leq C_HD_g$ then give, for an
acceptance $0\leq F\leq A$ with the stated path tests,
\begin{equation}\label{paper:ep-late-all-banks}
 e^{-C_v}\mathbb E_\pi[F\mid\operatorname{Class}_v]
       \leq C_He^{-b_0}W\prod_{g\ {\rm selected}}D_g.
\end{equation}
An early corner instead uses the literal bound
$q_v=\mathbb E_\pi[F\mid\operatorname{Class}_v]\leq1$ for
$0\leq F\leq1$.  It yields
$e^{-C_v}q_v\leq C_He^{-a_0}e^{-c_HV_l}$.
In the normalization $a_0=Y,b_0=0$, an independently verified same-family
target condition
\begin{equation}\label{paper:ep-target-contract}
 e^{-Y}\leq C_{\rm target}D_{\rm global},\qquad
 D_{\rm global}=W\prod_{g\ {\rm selected}}D_g,
\end{equation}
pays all early corners.  Combining with
\eqref{paper:ep-late-all-banks}, and the dimension-bounded number of
endpoint templates, gives for fixed $\delta>0$
\[
 \sum_v e^{-\delta C_v}q_v^\delta\leq C_\delta D_{\rm global}^\delta.
\]
Condition \eqref{paper:ep-target-contract} is an additional target
input, not a consequence of the count box.  In particular a short
prefix is never asked to pay count factors from an unrelated long
zero-rank tail.  The type envelope $A$ may be unbounded as long as the
actual early-corner acceptance $F$ is bounded by one.

\begin{theorem}[Intrinsic endpoint reserve]\label{paper:ep-intrinsic}
Let $V_{\rm old}$ be the total original old length and assume
\begin{equation}\label{paper:ep-incoming-size}
 a_0,b_0\geq0,\qquad b_0\leq B\sqrt{1+V_{\rm old}}.
\end{equation}
Define the intrinsic cost
\[
 K_v=Y-I_v=C_v-b_0\quad\hbox{when }C_v=a_0-I_v,
 \qquad Y=a_0-b_0.
\]
For every lawful endpoint template with $t_v=V_{l(v)}\geq V_{\rm old}$,
set
\begin{equation}\label{paper:ep-intrinsic-score}
 S_v^{\rm intr}=\begin{cases}
 h+A_{\rm in}+\overline C_v,&t_v\leq T/2,\\
 h+\overline C_v+S_{\rm good}(v),&t_v>T/2.
 \end{cases}
\end{equation}
For a fixed sufficiently large $h$, uniformly on the same count support,
\begin{equation}\label{paper:ep-intrinsic-comparison}
 c_HS_v^{\rm intr}\leq K_v+h\leq C_HS_v^{\rm intr},\qquad
 K_v\geq c_0\overline C_v-C_0.
\end{equation}
The assertion also holds for terminal $z=1$ using
Proposition~\ref{paper:ep-terminal-admission}; no proper-child
branching gap is needed for it.
\end{theorem}
\begin{proof}
If $V_{\rm old}\leq t_v\leq T/2$, the prefix reference charge
and \eqref{paper:ep-incoming-size} give, by one Young inequality,
$b_0\leq B\sqrt{1+t_v}\leq\overline C_v/4+C_B$.
Subtracting $b_0$ from \eqref{paper:ep-early-raw}, while retaining
$a_0$ with coefficient one, yields
\[
 K_v\geq a_0+\tfrac14\overline C_v-C_B.
\]
Its upper bound follows from the exact prefix formula and $b_0\geq0$:
$K_v\leq a_0+\overline C_v+C_H\sqrt{1+t_v}
\leq a_0+C_H'\overline C_v+C_H'$.
The case $t_v=V_{\rm old}=0$ is covered by the fixed constant.
For $t_v>T/2$, subtract $b_0$ in the exact tail formula itself:
\[
 K_v=Z_{l(v)}+\overline C_v+
                         \sum_{g\ {\rm proper}}\nu_{vg}.
\]
The proof of \eqref{paper:ep-late-raw} now gives directly
$K_v\geq\overline C_v/2+c_HS_{\rm good}-C_H$;
its upper proof gives
$K_v\leq C_H(\overline C_v+S_{\rm good})+C_H$.
Only the zero-tail and bad-group errors are absorbed, and full-group
errors cancel exactly.  These inequalities prove the theorem after
increasing $h$ once.
\end{proof}
On the zero-target packet, $b_0=\eta_{\rm prev}/u_{\rm prev}$
satisfies \eqref{paper:ep-incoming-size}: the preceding native length
is at most $V_{\rm old}$, its exit scale is bounded by a fixed multiple
of its square root, and $u_{\rm prev}$ is bounded below.  Terminal
root rounding adds only a fixed constant.  The complete-deadline
template has $K_v=0$ and intrinsic score $h$.  The all-zero template
has $K_v=a_0-b_0$ and is excluded when $V_{\rm old}>0$; that exclusion
is essential.  The late score in \eqref{paper:ep-intrinsic-score}
contains no $B_{\rm in}$ or preceding exit reserve.

\begin{corollary}[Legal native pairs]\label{paper:ep-intrinsic-pair}
For a genuine native synchronous edge in block $i=[l_i,b_i]$ having
at least one exterior cut, retain its lawful outside endpoint contexts.
Let $v^-,v^+$ denote their left and right completions and set
\begin{align*}
 X_i^{\rm intr}&=\min_vS_{v^-}^{\rm intr},&
 Y_i^{\rm intr}&=\min_vS_{v^+}^{\rm intr},\\
 x_i^{\rm intr}&=\min_v(K_{v^-}+h),&
 y_i^{\rm intr}&=\min_v(K_{v^+}+h).
\end{align*}
These actual minima are in the corresponding common scale boxes.
For a canonical full-to-deepest-$q$ edge they satisfy
\begin{equation}\label{paper:ep-intrinsic-increment}
 y_i^{\rm intr}-x_i^{\rm intr}
 =\Delta_i=(r_i-q)L_i-\log S_i\,n_i
                             +\log(q+z)K_{i,B}.
\end{equation}
A sufficient local pair, reading only this original local type, is
\begin{equation}\label{paper:ep-intrinsic-local}
 x_i^- =\max\{c_HX_i^{\rm intr},c_HY_i^{\rm intr}-\Delta_i\},
 \qquad y_i^-=x_i^-+\Delta_i.
\end{equation}
\end{corollary}
\begin{proof}
The actual name cut in $A\setminus B$ ensures every endpoint template
has $t_v\geq l_i\geq V_{\rm old}$, including a legal cut at the
old/native boundary.  Theorem~\ref{paper:ep-intrinsic} applies to
every member, so taking minima proves the common boxes.
Moving the synchronous cut across block $i$ changes only its original
rank and logarithmic count contribution.  This is exactly
\eqref{paper:ep-intrinsic-increment}, independent of outside context.
Subtracting the same $b_0$ from both raw endpoints changes no increment;
therefore the minimum right value equals the minimum left value plus
$\Delta_i$.  Lemma~\ref{paper:ep-local-pair} proves
\eqref{paper:ep-intrinsic-local}.  All original types, words and clocks
are retained.  Choose $h$ before any fixed reserve contraction so that
the contracted endpoints exceed its scalar threshold.
\end{proof}

\begin{proposition}[Allocation in a native-containing chamber]
\label{paper:ep-intrinsic-allocation}
In a lawful activation chamber containing a cut assigned to a native
physical block, the simultaneous cost allocation is valid with
$K=C-b_0$.  If $E_i$ is the original centered empirical bridge in
coordinate $i$, $L_{\rm intr}$ its synchronized multilinear endpoint
chord, and $\ell_i^{\rm intr}$ its legal outside-context minimum chord
shifted once by $C_0$, then, for $\theta=1/(2H)$,
\begin{equation}\label{paper:ep-intrinsic-allocation-equation}
 K_{\rm sync}+C_0
 \geq\tfrac12(L_{\rm intr}+C_0)
             +\sum_i[\theta\ell_i^{\rm intr}+E_i].
\end{equation}
For a surviving native canonical edge its sufficient support is
\begin{equation}\label{paper:ep-intrinsic-wall}
 \theta\ell_i^{\rm intr}(t)+E_i(t)
                           \geq g\,d_i(t)^\gamma-h_W,
\end{equation}
and the local pair \eqref{paper:ep-intrinsic-local} suffices for it.
\end{proposition}
\begin{proof}
Fix one native cut of the chamber.  Every endpoint vertex retains that
cut at a legal native boundary, hence has $t_v\geq V_{\rm old}$.
By Theorem~\ref{paper:ep-intrinsic}, every such vertex obeys
$K_v\geq c_0\overline C_v-C_0$ with $\overline C_v\geq0$.
The synchronization identity is unchanged on subtracting the constant
$b_0$ throughout this same-word chamber:
\[
 K_{\rm sync}=L_{\rm intr}+\sum_iE_i,
 \qquad L_{\rm intr}=L_{\rm multi}-b_0.
\]
The endpoint minima, shifted once, satisfy
$0\leq\ell_i^{\rm intr}\leq L_{\rm intr}+C_0$.
There are at most $H$ coordinates, so
$\theta\sum_i\ell_i^{\rm intr}\leq(L_{\rm intr}+C_0)/2$;
substitution proves \eqref{paper:ep-intrinsic-allocation-equation}.
There is no leftover $-\theta b_0$ term because the cost was normalized
before allocation.

Strong edges and reference-paid weak edges use the same named
regularity bounds and $K_v\geq c_0\overline C_v-C_0$.
For the remaining native canonical edges retain
\eqref{paper:ep-intrinsic-wall}.  An old coordinate uses its genuine
old charge (Proposition~\ref{paper:ep-old-weak}) and the same
regularity estimate.  For a proper atomic child, the branching
normal form and occupancy estimate then apply with this normalized
cost: synchronization differences, spread losses and pure-cell
lifetimes are unchanged by subtracting the chamber constant.
Divide the available reference, spread and curve budgets by the same
fixed fractions as in that allocation; every original cut and its
multiplicity remains.  At $z=1$ this argument first concerns star
histories, and extension to branching uses the separate floating-spine
argument with this repaired endpoint input.  Entirely old chambers
are outside the stated domain.
\end{proof}

A final two-color computation makes the normalization concrete.
On its last native block $[l,T]$, with every other fine name cut at
its own deadline, moving only the final boundary to $t$ gives the
exact global fine cost
\begin{equation}\label{paper:ep-two-color}
 K_{\rm fullstop}(t)=\log2\,N((t,T])-(T-t).
\end{equation}
Put $D=\log2\,N((l,T])-(T-l)$.  After a fixed shift its forward
pair is $(h+D,h)$ and its reverse pair $(h,h+D)$.
This context is among the legal outside contexts of the same $q=0$
coordinate.  Its shifted chord therefore dominates the repaired
minimum chord at both ends, and hence throughout the block.
Since $\ell_{\min}^{\rm intr}\geq0$, the sufficient allocated wall
implies
\[
 \theta\ell_{\min}^{\rm intr}+E\geq g d^\gamma-h_W
 \quad\Longrightarrow\quad K_{\rm fullstop}+h\geq g d^\gamma-h_W.
\]
More generally two lawful pairs at one exact count have the same
increment and differ only by a constant translation; their lower pair
is still lawful.  The ordinary all-reserve theorem evaluates the pair
in \eqref{paper:ep-two-color} only after verifying positivity, its
exact root count and unread-clock conditioning.  Up to fixed buffers
it has the form $(1,1+D)$, and contains no incoming $b_0$ term.
